\documentclass[11pt]{article}
\usepackage[margin=1in]{geometry}
\usepackage{amsmath,amssymb,booktabs}
\usepackage{setspace}
\onehalfspacing

\title{Tariffs and Exchange Rates in a Multilateral Setting: Results Outline}
\date{}

\begin{document}
\maketitle

\section{Overview}

The theoretical results are organized around a single question: how does the conventional two-country result that an import tariff appreciates the tariff-imposing country's currency transform once trade is multilateral? The progression is designed to separate three issues that coincide in a two-country model but become distinct when $N>2$: the adjustment of trade balances to exchange rates, the allocation of tariff-induced expenditure switching across foreign suppliers, and the distinction between bilateral and effective exchange-rate movements.

The analysis therefore begins with the general $N$-country equilibrium system, then uses the two-country model to recover the conventional Marshall-Lerner benchmark. A symmetric three-country model introduces the genuinely multilateral mechanism and delivers a transparent threshold for third-country exchange-rate reversal. The analysis then moves to general symmetric $N$, where the dependence of bilateral thresholds on the number of countries can be contrasted with the greater invariance of the nominal effective exchange rate (NEER). Asymmetry is introduced subsequently to determine which results are structural and which rely on symmetry. Alternative tariff configurations are considered last, since locally they can be understood as combinations of the primitive bilateral tariff shock.

\section{General $N$-Country Framework}

Begin with the general multilateral system. Although there are $N^2$ bilateral exchange-rate entries, no-arbitrage implies that bilateral exchange rates are generated by $N$ underlying currency values up to a common numeraire. There are therefore only $N-1$ independent relative exchange rates. Similarly, Walras' law implies that only $N-1$ of the $N$ trade-balance conditions are independent.

Choose country $A$ as the numeraire and collect the $N-1$ independent log exchange-rate changes in the vector $d\nu$. A local tariff perturbation can then be written as
\[
J\,d\nu=-F\,d\tau,
\]
or
\[
d\nu=(-J)^{-1}F\,d\tau.
\]

This representation separates two conceptually distinct objects. The Jacobian $J$ describes how the vector of trade balances responds to exchange-rate movements. The vector $F$ describes the impact of the tariff on trade balances at unchanged exchange rates, including expenditure switching, tariff-revenue effects, and income effects.

The role of $J$ is the natural multilateral extension of the adjustment mechanism underlying the Marshall-Lerner condition. In the two-country case, this object is scalar. In the multilateral case it is a matrix. Gross-substitution and M-matrix conditions provide useful sufficient conditions under which the exchange-rate adjustment system has a regular and predictable sign structure. Importantly, these conditions do not themselves determine the exchange-rate response to a tariff. The response also depends on the multidimensional tariff-incidence vector $F$.

This decomposition provides the organizing structure for the subsequent results:
\[
\text{tariff} \longrightarrow F
\longrightarrow \text{trade-balance disturbances}
\longrightarrow (-J)^{-1}
\longrightarrow \text{exchange-rate adjustment}.
\]

\section{The Two-Country Benchmark}

Specialize first to $N=2$. There is one independent trade balance and one independent exchange rate. Country $A$ imposes a tariff on imports from country $B$. At the initial exchange rate, the tariff reduces demand for imports from $B$ and improves $A$'s trade balance, after accounting for the associated tariff-revenue and income effects.

The equilibrium exchange-rate response satisfies
\[
\frac{de}{d\tau}
=
-\frac{\partial TB_A/\partial\tau}
{\partial TB_A/\partial e}.
\]
Under the Marshall-Lerner condition, the trade balance improves following a depreciation, so restoration of external balance following the tariff requires an appreciation of country $A$'s currency.

This establishes the conventional benchmark. It also highlights a restriction inherent in the two-country environment: substitution away from the tariffed foreign supplier is necessarily substitution either toward domestic goods or away from foreign expenditure altogether. There is no separate margin through which expenditure can be redirected toward an untreated foreign supplier.

\section{Three Countries: The Multilateral Mechanism}

\subsection{Symmetric benchmark}

Introduce a third country $C$ while initially maintaining symmetry. Country $A$ imposes a tariff on imports from $B$, while imports from $C$ remain untariffed. Preferences distinguish two substitution margins. The elasticity $\eta$ governs substitution between domestic tradables and the foreign tradable composite, while $\rho$ governs substitution across foreign origins.

The tariff-incidence object $F$ is now multidimensional. A tariff on $B$ can redirect expenditure not only toward domestic production but also toward $C$. The exchange-rate adjustment system can continue to satisfy the multilateral analogue of the Marshall-Lerner condition, yet the tariff-induced disturbances need not point in the same direction across countries.

The symmetric three-country solution yields two qualitatively different bilateral responses. Country $A$ appreciates against the tariff target $B$, while the response against the untreated country $C$ is ambiguous. There exists a threshold $\rho_3^*$ such that
\[
\frac{d\log e_{AC}}{d\tau}>0
\quad\Longleftrightarrow\quad
\rho>\rho_3^*,
\]
under the regularity condition ensuring the normal sign of the denominator.

The threshold result isolates the genuinely multilateral mechanism. A sufficiently high cross-origin substitution elasticity makes expenditure diversion toward $C$ strong enough that equilibrium requires $A$ to depreciate against $C$, even while $A$ appreciates against $B$. The failure of the conventional bilateral prediction therefore does not require an abnormal exchange-rate adjustment mechanism. It can arise entirely from the multidimensional structure of the tariff-induced trade-balance disturbance.

\subsection{Bilateral versus effective exchange rates}

The three-country setting also introduces a distinction that does not exist for $N=2$: ``the exchange rate'' is no longer a unique scalar object. Bilateral exchange rates can move in opposite directions, while an effective exchange rate summarizes the aggregate external value of the currency.

For symmetric trade weights, define country $A$'s log effective exchange rate as the average of its bilateral log exchange rates. In the unilateral-tariff experiment, the terms involving the cross-origin substitution elasticity cancel from the numerator of the NEER response. The resulting expression takes the form
\[
\frac{d\nu_A^{NEER}}{d\tau}
=
-\frac{\rho_1^*}{D_N},
\]
where
\[
\rho_1^*
=
1+\alpha_D(\eta-1)-\alpha_T(1-\alpha_D)
=
\alpha_D\eta+(1-\alpha_D)(1-\alpha_T).
\]
Under economically admissible parameter values, $\rho_1^*>0$, and under the standard regularity condition $D_N>0$. The unilateral tariff therefore appreciates country $A$'s NEER even when it depreciates against an untreated trading partner.

This result suggests a useful interpretation of the multilateral adjustment. The exchange-rate vector contains an aggregate home-versus-foreign component and a relative reallocation component across foreign suppliers. Cross-origin substitution can strongly affect the latter, and can reverse individual bilateral responses, without overturning the sign of the aggregate effective exchange-rate response in the symmetric benchmark.

\section{General Symmetric $N$}

The next step is to ask how the multilateral mechanism scales as the number of trading partners increases. Maintain symmetry and continue to consider the primitive experiment in which $A$ imposes a tariff on one country $B$.

Define
\[
D_N
=
N\alpha_D(\eta-1)+(N-2)\rho+1.
\]
The bilateral response against an untreated country $C$ can be expressed using the threshold
\[
\rho_N^*=N\rho_1^*.
\]
Thus,
\[
\frac{d\log e_{AC}}{d\tau}>0
\quad\Longleftrightarrow\quad
\rho>N\rho_1^*.
\]

The bilateral reversal threshold therefore rises with $N$. In the symmetric experiment, total foreign expenditure is divided across a larger number of suppliers, so diversion away from $B$ is spread across more untreated countries. A larger cross-origin substitution elasticity is consequently required for the response against any particular $C$ to reverse.

The effective exchange-rate result is markedly different. With equal weights across foreign partners,
\[
\frac{d\nu_A^{NEER}}{d\tau}
=
-\frac{\rho_1^*}{D_N}.
\]
The number of countries and the cross-origin elasticity affect the magnitude of the response through $D_N$, but not its sign. The bilateral consequences of the tariff therefore become increasingly dependent on the dimensionality of the trading system, while the conventional effective appreciation result survives aggregation in the symmetric unilateral-tariff benchmark.

This contrast is central to the interpretation of the model. The dependence of $\rho_N^*$ on $N$ primarily describes how the adjustment is distributed across individual trading partners. Aggregating those partner-specific movements into an effective exchange rate removes the corresponding dilution effect from the sign condition.

\section{Asymmetry: What Is Structural?}

Symmetry is useful for isolating the mechanism and obtaining transparent thresholds, but the next question is which results survive when bilateral trade shares differ.

In the asymmetric three-country model, the same economic channels remain present, but the threshold for third-country depreciation becomes a function of bilateral expenditure shares and country exposures. The purpose of this extension is therefore not to introduce a new mechanism, but to show how the common multilateral mechanism is reweighted by the structure of the trade network. The symmetric threshold is recovered as a special case.

Asymmetry is potentially more consequential for the effective exchange rate. The exact cancellation of the cross-origin reallocation terms in the symmetric NEER need not hold under arbitrary trade weights. It is therefore useful to discipline asymmetry in a way that remains comparable as $N$ changes.

One natural parameterization lets the tariff target $B$ have a fixed multiple $\kappa$ of the symmetric bilateral foreign share:
\[
\beta_{AB}
=
\frac{\kappa}{N-1},
\]
with the remaining foreign expenditure distributed across the $N-2$ untreated countries. Here $\kappa=1$ corresponds to exact symmetry, while $\kappa>1$ or $\kappa<1$ makes $B$ respectively more or less important than a typical foreign partner.

Under this scaling, asymmetry remains economically comparable as $N$ changes rather than making $B$ increasingly exceptional. The large-$N$ NEER response takes the approximate form
\[
\frac{d\nu_A^{NEER}}{d\tau}
\approx
-\frac{\kappa}{N-2}
\frac{\rho_1^*}
{\rho+\alpha_D(\eta-1)}.
\]
Under normal parameter values, the response therefore approaches zero from the appreciation side. Exact symmetry gives effective appreciation for every $N$, while proportionally scaled asymmetry preserves the conventional sign asymptotically. Finite-$N$ asymmetry can nevertheless weaken or, for sufficiently uneven trade exposures, overturn the aggregate result.

The role of the asymmetric analysis is therefore to distinguish the robustness of the aggregate result from the much greater sensitivity of bilateral responses. It also clarifies that unrestricted asymmetry is not an informative generalization across $N$ unless the degree of trade concentration is itself disciplined.

\section{Alternative Tariff Configurations}

Only after establishing the primitive bilateral tariff response is it necessary to vary the tariff experiment. Since the analysis concerns local first-order responses, more complicated tariff configurations can be constructed as linear combinations of directed bilateral tariff shocks. These experiments are therefore best viewed as applications of the preceding theory rather than as separate models.

\subsection{Broad unilateral protection}

Consider first a tariff imposed by $A$ symmetrically on all foreign suppliers. This experiment removes the discriminatory cross-origin margin that drives the third-country reversal in the bilateral tariff case. It therefore provides a useful benchmark for distinguishing the aggregate home-versus-foreign effect of protection from the redistribution of expenditure across foreign origins. The precise exchange-rate result should be derived from the general $N$-country system, but the experiment is conceptually useful because it shuts down the central diversion mechanism by symmetry.

\subsection{Symmetric bilateral trade war}

A symmetric trade war between $A$ and $B$ can be constructed by adding the response to an $A$-on-$B$ tariff and the mirror-image response to a $B$-on-$A$ tariff. By symmetry, the direct $A$-$B$ exchange-rate movements cancel, while the effects relative to untreated countries reinforce.

The resulting third-country and effective exchange-rate response is governed by the cross-origin substitution margin. Under the symmetric benchmark,
\[
\frac{d\nu_A^{NEER}}{d\tau}>0
\quad\Longleftrightarrow\quad
\rho>\rho_1^*.
\]
Thus, unlike the unilateral bilateral tariff, reciprocal protection can cause both tariffing countries to depreciate effectively against the rest of the world when cross-origin substitution is sufficiently strong.

The trade-war result is an important qualification to the NEER preservation result, but it does not require a new mechanism. Reciprocity changes the combination of primitive tariff shocks: it cancels the bilateral appreciation component between the belligerents and allows the third-country diversion component to determine the aggregate effective exchange-rate response.

\section{Synthesis}

The progression from two to many countries changes the interpretation of the conventional tariff-exchange-rate result. In a two-country model, the bilateral exchange rate, the effective exchange rate, and the allocation of expenditure across foreign suppliers all collapse into a single margin. A tariff therefore produces a simple appreciation result under the Marshall-Lerner condition.

In a multilateral system, these margins separate. The exchange-rate adjustment mechanism can remain conventional, while the tariff generates a multidimensional vector of trade-balance disturbances. Cross-origin substitution then determines how adjustment is distributed across trading partners and can reverse the exchange-rate response against untreated countries. The corresponding bilateral threshold depends on the number and structure of trading partners.

At the same time, the symmetric unilateral-tariff benchmark preserves the conventional result at the effective-exchange-rate level. This provides a precise sense in which the two-country conventional wisdom survives multilateralization: it is considerably more robust as a statement about the aggregate external value of the tariffing country's currency than as a statement about every bilateral exchange rate.

Asymmetry determines how cleanly the aggregate and relative components separate, while alternative tariff configurations determine how the primitive bilateral responses are combined. In particular, reciprocal tariffs can cancel the conventional bilateral appreciation component and allow the multilateral diversion channel to determine even the effective exchange-rate response.

The main theoretical message is therefore not simply that tariffs can appreciate or depreciate currencies. Rather, it is that moving from a bilateral to a multilateral trading system changes the dimensionality of the exchange-rate adjustment problem. The conventional result can survive in aggregate while failing systematically at the bilateral level, and the model characterizes the substitution, network, and policy conditions under which each component dominates.

\end{document}
